% Author: JoonHo Lee (jlee296@ua.edu)
\begin{table}[tbp]\centering
\caption{Score-test rejection when latent means are equal}
\label{tab:main-null}
{\small\begin{tabular}{llrrr}
\toprule
Focal SD & Score & Index gap & Rejection & MCSE \\
\midrule
$1.0$ & Sum & $-0.0004$ & $0.035$ & $0.013$ \\
$1.0$ & EAP & $-0.0004$ & $0.040$ & $0.014$ \\
$1.0$ & WLE & $-0.0004$ & $0.045$ & $0.015$ \\
\addlinespace
$0.6$ & Sum & $0.187$ & $0.040$ & $0.014$ \\
$0.6$ & EAP & $0.187$ & $0.035$ & $0.013$ \\
$0.6$ & WLE & $0.187$ & $0.045$ & $0.015$ \\
\addlinespace
\bottomrule
\end{tabular}}
\floatnote{Each condition uses one 20-item 2PL form, held fixed for 200 replications with 500 persons per group. Reference ability is $N(0,1)$; focal ability has SD 1 or .6. The index gap is reference minus focal. Tests compare score means using two-sided Welch tests at .05. Equal latent means need not imply equal score means when dispersions differ. Monte Carlo standard errors exclude variation across forms.}
\end{table}
